\documentclass{report}
\usepackage{amsmath,amssymb}
\setlength{\parindent}{0mm}
\setlength{\parskip}{1em}
\begin{document}
\begin{center}
\rule{6in}{1pt} \
{\large William, T Taitano \\
{\bf Development of a fully implicit physics-based multi-grid Accelerator for Vlasov-Ampere Particle in Cell (PIC) system}}

5577 Wyoming Blvd NE \\ 207N \\ Albuquerque \\ NM \\ 87109
\\
{\tt williamtaitano1208@hotmail.com}\\
Dana, A Knoll\\
Luis Chacon\end{center}

\noindent Explicit time integration is predominant in particle-in-cell
(PIC) methods for kinetic plasma
simulation. In explicit methods, the time-step size is restricted by the
CFL condition. This makes
detailed investigation of problems that occurs on ion time-scales
challenging. Implicit methods
are not restricted by any of the numerical time-scales.
Jacobian-Free-Newton-Krylov (JFNK) is one method that solves a coupled
set of non-linear system of equations that promises investigation of
implicit multi-time-scale kinetic plasma physics simulation possible [1].
With appropriate discretization, one can also use a mesh spacing which is
larger than the Debye length. An alternative implicit approach have also
been investigated by Degond et al [2]. The approach is formulated around
a semi-implicit, reformulated Poisson's equation in order to remove the
stiffness associated with the pure Poisson equation to solve for the
electric-field. \\

\noindent Our work is based on developing a moment-based approach which
may allow us to abandon the tight JFNK iteration between fields and
particles. The work is an extension of [1]. The method achieves
acceleration of the fine-scale transport equation solution that lives in
a 6D phase-space through the coarse-resolution moment system which lives
in the 3D physical space. This is analogous to the multi-grid method in
which the coarse grid solution damps out the short wavelength solutions
and accelerates the solution of the fine grid system through series of
prolongation and restriction steps for iterative methods. In our method,
we will perform a similar prolongation and restriction process between
the two physical scale. The prolongation process that transfers
information from the coarse scale moment system to the fine scale
transport system can be regarded as the coarse space moment expansion of
transport operators (such as the scattering operator in neutron
transport). The restriction process, which transfers information from the
fine-scale transport system to the coarse-scale moment system is the
phase-space averaging of the transport solution to obtain quantities for
the moment system (such as the Eddington tensor in neutron transport or
the stress tensor in plasma physics). \\

\noindent Hence, the approach may be truly interpreted as a
"physics-based" multi-grid method. We present our progress on this new
multi-scale algorithm and contrast it with the approaches of [1] and [2].
We will also show a connection to moment-based acceleration of transport
iteration used to solve for the neutron transport equation.\\[4pt] [1] G.
Chen, et al., J. Comp. Phys., vol. 230 (2011)
\\[0pt] [2] P. Degond, et al., J. Comp. Phys., vol. 229 (2010)


\end{document}
